\section{A relative derived centre and its nonpositive part}
\label{v4-arith:ks:relative-centre}

A strict associative differential graded algebra can retain higher
multiplications on its cohomology. Take $\mathbb Z[t]\langle\tau\rangle$, with central $t$ and
\[
 |t|=0,\qquad |\tau|=-1,\qquad d\tau=t^2.
\]
The relation $[t]^2=0$ has primitive $\tau$, whose square
$u=\tau^2$ is a nonzero class of degree $-2$.
The fourth operation will record this square.
The same calculation occurs inside a relative derived centre, where
an additional class of degree two changes the multiplication relation.

Fix a commutative unital ring $S$ and an integer $r\geq1$.
All differentials have cohomological degree one.
Algebras are unital and associative, with a central closed $S$-action.
Tensor products and polynomial sums are algebraic.
A homogeneous tensor of maps acts by
\[
 (f\otimes g)(x\otimes y)=(-1)^{|g||x|}f(x)\otimes g(y).
\]
An odd generator need not have square zero in this category.
Formality over $S$ means a finite zigzag of unital $S$-linear
DGA quasi-isomorphisms to cohomology with zero differential.
No additional generator is fixed by this definition.

Put $R=S[t]$ and
\[
 E_r=R\langle\tau\rangle,\qquad d\tau=t^r.
\]
The variable $t$ commutes with $\tau$.
Over the characteristic-zero geometric base, this is the algebra in
Booth, \emph{The derived contraction algebra}, Lemma~4.3.8, p.~22,
with $(\xi,\zeta,n)=(t,\tau,r)$.
The following calculation uses only its displayed algebraic presentation.

The relative derived centre is
\[
 Z_R(E_r)=\operatorname{RHom}_{E_r\otimes_R E_r^{\mathrm{op}}}
 (E_r,E_r),
\]
with its derived composition product.
The opposite algebra uses the Koszul sign convention.
The word ``relative'' fixes $R=S[t]$ in the bimodule tensor product.
In particular, this construction differs from $Z_S(E_r)$.

\begin{proposition}[a multiplicative model for the relative centre]
\label{v4-arith:ks:centre-model}
\ClaimStatusProvedHere{}
The relative centre has the associative DGA model
\[
 C_{S,r}=S[t]\langle\tau\rangle[z]/(z^2),\qquad
 |z|=2,\qquad
 D\tau=t^r-2\tau^2z,\quad Dt=Dz=0,
\]
where $z$ commutes with $t,\tau$.
There is a multiplicative quasi-isomorphism from this algebra to the
endomorphism DGA of a free bimodule resolution of $E_r$.
\end{proposition}

\begin{proof}
Let $P$ be the free graded $E_r$-bimodule on generators
$e_0,e_1$ of degrees $0,-2$, with
\[
 d_Pe_0=0,\qquad d_Pe_1=\tau e_0-e_0\tau.
\]
The Leibniz rule defines the differential on all of $P$.
Its square on $e_1$ is $t^re_0-e_0t^r=0$.
Define $\epsilon:P\to E_r$ by
$\epsilon(e_0)=1$ and $\epsilon(e_1)=0$.

The tensor-algebra sequence
\[
 0\longrightarrow E_r\otimes_R R\tau\otimes_R E_r
 \longrightarrow E_r\otimes_R E_r
 \xrightarrow{\mathrm{mult}} E_r\longrightarrow0
\]
is exact as a sequence of graded modules.
The first map sends $a\otimes\tau\otimes b$ to
$a\tau\otimes b-a\otimes\tau b$.
At each total word length, its matrix is the difference map between
successive cuts of one word.
Its image consists exactly of combinations whose coefficients sum to zero,
and its kernel vanishes by inspecting the first cut, then each successive cut.
The universal differential of $d\tau=t^r$ relative to $R$ is zero.
Thus this sequence is a sequence of complexes.
Its two-term totalization, with the usual shift signs, is $P$.
Exactness proves that $\epsilon$ is a quasi-isomorphism.
The filtration generated first by $e_0$ and then by $e_1$ makes $P$
a semifree bimodule, so it computes the displayed derived Hom.

There is a counital coassociative diagonal
$\Delta:P\to P\otimes_{E_r}P$ given by
\[
 \Delta(e_0)=e_0\otimes e_0,\qquad
 \Delta(e_1)=e_1\otimes e_0+e_0\otimes e_1.
\]
The two middle terms in $d\Delta(e_1)$ cancel in the balanced tensor product.
This proves that $\Delta$ commutes with the differential.
Coassociativity and the counit identities follow on these two generators.

A homogeneous bimodule map of degree $k$ satisfies
$f(apb)=(-1)^{k|a|}af(p)b$.
Identify $\operatorname{Hom}_{E_r^e}(P,E_r)$ with $E_r\oplus E_rz$
by the values on $e_0,e_1$.
Convolution through $\Delta$ gives
\[
 (a+bz)(a'+b'z)=aa'+(ba'+ab')z.
\]
The differential on the value $a$ has additional $e_1$-component
\[
 -\tau a+(-1)^{|a|}a\tau.
\]
For $a=\tau$, this is $-2\tau^2$.
These are exactly the multiplication and differential of $C_{S,r}$.
In particular $D(\tau^2)=0$, so $D^2=0$.

For a convolution cochain $f$, the map
\[
 \rho(f)=(1\otimes f)\Delta:P\longrightarrow P
\]
uses the right action to identify $P\otimes_{E_r}E_r$ with $P$.
Coassociativity gives $\rho(fg)=\rho(f)\rho(g)$, with the tensor signs above.
The counit gives $\epsilon\rho(f)=f$.
Since $P$ is semifree, composition with $\epsilon$ is a
quasi-isomorphism from $\operatorname{End}_{E_r^e}(P)$ to
$\operatorname{Hom}_{E_r^e}(P,E_r)$ as complexes.
Its section $\rho$ is therefore the claimed DGA quasi-isomorphism.
\end{proof}

To separate the coefficient in the differential from the base ring,
let $c\in S$ and define
\begin{equation}\label{v4-arith:ks:coefficient-family}
 \begin{gathered}
 B=S[u,v]/(v^2),\quad T=B[t],\quad f=t^r-cv,\\
 A(S,c,r)=T\oplus\tau T,\qquad
 |u|=-2,\quad |v|=|t|=0,\quad |\tau|=-1,\\
 \tau^2=u,\qquad d(a+\tau b)=fb.
 \end{gathered}
\end{equation}
Here $\tau$ commutes with $T$.
The equality $d(\tau^2-u)=f\tau-\tau f=0$ verifies the defining relation.
This family allows arbitrary $c$.
The centre in Proposition~\ref{v4-arith:ks:centre-model} has the specific coefficient $c=2$.

\begin{proposition}[the truncation map and its quotient]
\label{v4-arith:ks:centre-truncation}
\ClaimStatusProvedHere{}
There is an injective DGA map
\[
 j:A(S,2,r)\longrightarrow C_{S,r},\qquad
 j(t)=t,\quad j(\tau)=\tau,\quad j(u)=\tau^2,\quad j(v)=\tau^2z.
\]
Its image is the nonpositive truncation $\tau_{\leq0}C_{S,r}$.
It induces an isomorphism on cohomology in every nonpositive degree.
The quotient complex has its only cohomology in degree two, where it is
$S[t]/(t^r)\,z$.
\end{proposition}

\begin{proof}
Write $u=\tau^2$ inside $C_{S,r}$.
Its even and odd parts are
\[
 S[u,z,t]/(z^2)\quad\hbox{and}\quad
 \tau S[u,z,t]/(z^2).
\]
The monomials $u^at^i$ and $u^azt^i$ are $S$-linearly independent.
Sending $u^avt^i$ to $u^{a+1}zt^i$ proves injectivity of $j$,
also on the odd summand.
The map respects the differential because $t^r-2v$ maps to $D\tau$.
All degree-zero elements of $C_{S,r}$ are closed.
Its only positive terms are $S[t]\tau z$ in degree one and $S[t]z$ in degree two.
The quotient of $j$ is therefore the complex
\[
 S[t]\tau z\xrightarrow{\ t^r\ }S[t]z
\]
in degrees one and two.
Multiplication by $t^r$ is injective, which proves the assertions.
\end{proof}

\begin{remark}[the missing positive generator]
For $S\ne0$, the map $j$ is not a quasi-isomorphism.
It has no graded algebra retraction: every such retraction sends $z$ to zero,
but a retraction would also send $uz$ to the nonzero element $v$.
The multiplicative quasi-isomorphism of Proposition~\ref{v4-arith:ks:centre-model}
identifies the full $C_{S,r}$ with the relative centre.
The map $j$ identifies $A(S,2,r)$ with its nonpositive truncation.
\end{remark}

\section{The fourth operation over a coefficient ring}
\label{v4-arith:ks:fourth-operation}

The relation $t^r=cv$ has a chosen primitive $\tau$.
Products of two such primitives produce $u$.
To decide whether changing the primitives can remove this effect,
we first compute the transferred operation and then its cohomology class.

Multiplication by the monic polynomial $f=t^r-cv$ on $T$ is injective.
A nonzero polynomial times $f$ retains its nonzero leading coefficient,
including when $S$ has zero divisors.
It follows that
\begin{equation}\label{v4-arith:ks:family-cohomology}
 H=H^*(A(S,c,r))=S[u,v,t]/(v^2,t^r-cv),\qquad H^{\mathrm{odd}}=0.
\end{equation}
Division by the monic polynomial supplies the $S$-basis
\[
 u^av^jt^i,\qquad a\geq0,\quad j\in\{0,1\},\quad 0\leq i<r.
\]
Each $H^{-2a}$ is free of rank $2r$.

Let $i:H\to T$ choose the representative of $t$-degree less than $r$.
Let $p:A\to H$ take the even component modulo $f$ and kill the odd component.
Set
\[
 h(x)=\tau\frac{x-ip(x)}f\quad(x\in T),\qquad h(\tau T)=0.
\]
Existence and uniqueness of monic division give
\[
 dh+hd=1-ip,\qquad pi=1,\qquad h^2=hi=ph=0.
\]
Define the $B$-bilinear degree-zero map $F:H^{\otimes_S2}\to H$ by
\[
 i(x)i(y)-i(xy)=f\,iF(x,y).
\]
On the normal basis it is
\[
 F(t^i,t^j)=
 \begin{cases}
 t^{i+j-r},&i+j\geq r,\\
 0,&i+j<r.
 \end{cases}
\]
For the first case, use
$t^{i+j}=f t^{i+j-r}+cv t^{i+j-r}$.
Both quotients have degree less than $r$.

We use $A_\infty$ operations of degree $|m_n|=2-n$
and morphism components of degree $|f_n|=1-n$.
Explicitly, the operation identities have signs
\[
 \sum_{a+b+q=n}(-1)^{a+bq}
 m_{a+1+q}(1^{\otimes a}\otimes m_b\otimes1^{\otimes q})=0.
\]
The indices satisfy $a,q\geq0$ and $b\geq1$.
In a morphism identity, the term
$m_k(f_{i_1}\otimes\cdots\otimes f_{i_k})$
has sign $(-1)^{\sum_{j=1}^k(k-j)(i_j-1)}$.
Tensor evaluation also uses the Koszul rule fixed above.
A minimal structure has $m_1=0$.

\begin{lemma}[transfer and the first obstruction]
\label{v4-arith:ks:transfer-four}
\ClaimStatusProvedHere{}
The contraction gives a minimal structure on $H$ and a morphism to $A$
whose first components are
\[
 f_1=i,\qquad f_2=-\tau iF,\qquad f_3=m_3=0,
\]
\begin{equation}\label{v4-arith:ks:m-four}
 m_4(x,y,z,w)=-uF(x,y)F(z,w),\qquad
 f_4(x,y,z,w)=\tau u\,iF(F(x,y),F(z,w)).
\end{equation}
The operation $m_4$ is a Hochschild cocycle of bidegree $(4,-2)$.
\end{lemma}

\begin{proof}
For completeness, the transfer can be constructed over $S$ without a field hypothesis.
Set
\[
 U_n=\sum_{a+b=n}(-1)^{a-1}\mu(f_a\otimes f_b),\qquad
 m_n=pU_n,\qquad f_n=-hU_n\quad(n\geq2),
\]
where $a,b\geq1$ and $\mu$ is multiplication in $A$.
Each formula uses only smaller components.

Here is the compatibility argument.
Suspend by $|sx|=|x|-1$, and use the tensor coalgebra
$T^c(sH)=\bigoplus_{n\geq1}(sH)^{\otimes_S n}$ with deconcatenation.
Write $I,P,K$ for the suspended maps $i,p,h$.
Write $D_1$ for the suspended differential of $A$ and $\Lambda$
for its binary multiplication coderivation, with the signs just specified.
Thus $D_1^2=\Lambda^2=D_1\Lambda+\Lambda D_1=0$.
The recursion becomes
\[
 Q_n=\sum_{a+b=n}\Lambda_2(\Phi_a\otimes\Phi_b),\quad
 \Phi_1=I,\quad \Phi_n=-KQ_n,\quad \beta_1=0,\quad\beta_n=PQ_n.
\]
Extend $\Phi$ by consecutive partitions and $\beta$ by insertion.
These give a coalgebra map and a coderivation respectively.

Let $E=(D_1+\Lambda)\Phi-\Phi\beta$.
If $E$ vanishes below length $n$, the coproduct identity
\[
 \Delta E=(E\otimes\Phi+\Phi\otimes E)\Delta
\]
makes $E$ primitive in length $n$.
A primitive element has length one, since the cut after its first factor
detects every component of greater length.
Writing $Q$ for the corestriction of $\Lambda\Phi$, one obtains
\[
 0=\pi\Lambda E=-D_1Q-Q\beta,
 \qquad
 \pi E=K(D_1Q+Q\beta)=0.
\]
The second identity uses $D_1K+KD_1=1-IP$.
Induction gives $E=0$.
The highest-length component of $\Phi$ is $I^{\otimes n}$,
which has left inverse $P^{\otimes n}$.
Thus $\Phi$ is injective on finite tensor sums, and
$\Phi\beta^2=(D_1+\Lambda)^2\Phi=0$ implies $\beta^2=0$.
This proves the transfer identities in every arity.

The maps $p,h$ kill odd outputs, giving $f_3=m_3=0$.
Only the partition $(2,2)$ contributes to $m_4$, with a negative sign.
Its product is $-u\,iF(x,y)iF(z,w)$.
Applying $p$ and $-h$ gives \eqref{v4-arith:ks:m-four}.
For an even algebra the Hochschild differential is
\[
 \begin{split}
 (\delta K)(x_1,\ldots,x_{n+1})={}&x_1K(x_2,\ldots,x_{n+1})\\
 &+\sum_{j=1}^n(-1)^jK(x_1,\ldots,x_jx_{j+1},\ldots,x_{n+1})\\
 &+(-1)^{n+1}K(x_1,\ldots,x_n)x_{n+1}.
 \end{split}
\]
The arity-five identity gives $\delta m_4=0$.
Alternatively, expand the representative of a triple product in both
parenthesizations, cancel the injective factor $f$, and apply $p$.
This gives $\delta F=0$, hence $\delta(F\smile F)=0$.
\end{proof}

A nonzero value of $m_4$ alone does not prove nonformality.
A change of coordinates can alter it by a Hochschild boundary.
An arbitrary formality zigzag must therefore be compared with the
particular minimal structure just constructed.

\begin{lemma}[finite comparison over an arbitrary ring]
\label{v4-arith:ks:projective-transport}
\ClaimStatusProvedHere{}
Suppose a DGA $D$ has an explicit minimal model $V$ whose homogeneous
components are projective $S$-modules.
That model can be transported through any finite quasi-isomorphism zigzag,
through every specified finite arity.
Consequently, formality of $A(S,c,r)$ implies
\[
 m_4=\delta K
\]
for an $S$-linear cochain $K:H^{\otimes_S3}\to H$ of degree $-2$.
The cochain need not preserve $t,u,v$.
\end{lemma}

\begin{proof}
First take a degreewise surjective DGA quasi-isomorphism $q:P\to D$.
Its kernel $L$ is acyclic.
Suppose a coalgebra morphism from the minimal source has been lifted
through arity $n-1$.
Every degree of $W_n=(sV)^{\otimes_S n}$ is projective:
tensor products of projective modules are projective, and these degrees
are direct sums of such tensor products.
Choose a graded lift of the arity-$n$ component through $q$.

The defect $E=D_P\Phi-\Phi b$ is primitive in arity $n$ because
all smaller defects vanish.
It lies in $sL$ because its image under $q$ vanishes.
The identity $D_PE+Eb=0$ makes it closed:
minimality means that $b$ lowers tensor length, so $Eb$ vanishes here.
Acyclicity supplies the surjection
\[
 d:(sL)^j\longrightarrow Z^{j+1}(sL).
\]
Projectivity of $W_n^j$ lifts this particular defect through $d$.
Subtract the lift from the chosen morphism component.
This removes the defect without changing its image under $q$.
No contraction of the acyclic module $L$ is required.

Every backward arrow admits such a surjective replacement.
Let $J$ be the free $S$-module on $e_0,e_1,\eta$, with
$|e_i|=0$, $|\eta|=1$,
\[
 e_i^2=e_i,\quad e_0e_1=e_1e_0=0,\quad
 e_0\eta=\eta e_1=\eta,\quad
 \eta e_0=e_1\eta=\eta^2=0,
\]
\[
 de_0=-\eta,\qquad de_1=\eta,\qquad d\eta=0.
\]
The unit is $e_0+e_1$.
These are upper triangular matrix multiplication and its differential.
The endpoint maps $\mathrm{ev}_i:J\to S$ select the diagonal coefficients.
The map $k(\eta)=e_1$, $k(e_i)=0$ satisfies
\[
 dk+kd=1-\mathrm{const}\,\mathrm{ev}_0.
\]
On $D\otimes_S J$, use $(-1)^{|a|}a\otimes k(\omega)$ on $a\otimes\omega$.
This gives the same contraction with the tensor differential.

For a quasi-isomorphism $g:D'\to D$, form the pullback DGA
\[
 P(g)=\{(b,\omega)\in D'\times(D\otimes_S J):
             \mathrm{ev}_0\omega=g(b)\}.
\]
Projection $p:P(g)\to D'$ is a homotopy equivalence as a complex,
with section $j(b)=(b,g(b)\otimes1)$ and the displayed contraction.
The other endpoint $q=\mathrm{ev}_1$ is degreewise surjective:
$(0,a\otimes e_1)$ lifts every homogeneous $a\in D$.
Since $qj=g$, it is also a quasi-isomorphism.
Thus each backward arrow can be replaced by the span
$D'\xleftarrow{p}P(g)\xrightarrow{q}D$.
Lift through $q$ and compose with $p$.
Forward arrows require only composition.
This traverses the finite zigzag while keeping the minimal source fixed.

At a strict cohomology target, the linear component is a graded algebra automorphism.
Compose with its inverse to make that component the identity.
Evenness of $H$ forces the second component to vanish.
The arity-four morphism equation now reads $m_4=\delta K$,
with $K$ the third component.
\end{proof}

\begin{theorem}[the coefficient criterion for formality]
\label{v4-arith:ks:coefficient-formality}
\ClaimStatusProvedHere{}
For every commutative unital ring $S$, every $c\in S$, and every $r\geq1$,
\[
 \boxed{\quad A(S,c,r)\text{ is formal over }S
       \quad\Longleftrightarrow\quad r=1\ \text{or}\ c\in S^\times.\quad}
\]
\end{theorem}

\begin{proof}[Necessity]
Assume $r\geq2$ and put $a=t^{r-1}$.
The three adjacent products in $(a,t,a,t)$ are $cv$.
Since $F(a,t)=1$, one has
\[
 m_4(a,t,a,t)=-u.
\]
Evaluation $t,v\mapsto0$ defines a graded algebra map $H\to S[u]$.
Follow it in degree $-2$ by extraction of the coefficient of $u$.
The resulting functional $\lambda:H^{-2}\to S$ satisfies
$\lambda(u)=1$ and $\lambda(tH^{-2})=0$.
For every $S$-linear degree-$-2$ ternary cochain,
\[
 \begin{split}
 (\delta K)(a,t,a,t)={}&aK(t,a,t)+K(a,t,a)t\\
 &+c\bigl(-K(v,a,t)+K(a,v,t)-K(a,t,v)\bigr).
 \end{split}
\]
Applying $\lambda$ kills the exterior terms because $a\in tH^0$.
Thus $\lambda(\delta K(a,t,a,t))\in cS$.
The comparison lemma makes formality imply $-1\in cS$, so $c$ is a unit.
Equivalently, evaluation by $\lambda$ modulo $cS$ kills every boundary
and sends $m_4$ to $-1$ in $S/cS$.
This argument reduces one cochain equation modulo $cS$.
It does not tensor an arbitrary quasi-isomorphism with a nonflat quotient.
\end{proof}

\section{A multiplicative resolution in the formal case}
\label{v4-arith:ks:positive-resolution}

For $r=1$, the inclusion $B\hookrightarrow A(S,c,1)$ is a quasi-isomorphism,
since its cohomology is $B[t]/(t-cv)\cong B$.
For unit $c$ and arbitrary $r$, eliminating $v=c^{-1}t^r$ gives
\[
 H\cong S[t,u]/(t^{2r}).
\]
A resolution must account for the square of the primitive of $t^{2r}$.

Put $R'=S[t,u]$ and $s=t^r$.
The only element of degree $-1$ in $A$ whose differential is $s^2$ is
$\tau(s+cv)$, by injectivity of $d$ on the odd summand.
Its square is
\[
 u(s^2+2csv).
\]
For $S\ne0$, this polynomial is nonzero.
Thus the exterior resolution with generator $h$, $h^2=0$, $dh=s^2$
has no DGA map to $A$ over $R'$.
Successive generators remove these successive multiplication defects.

\begin{proposition}[explicit quasi-isomorphisms]
\label{v4-arith:ks:positive-zigzag}
\ClaimStatusProvedHere{}
Let
\[
 P=R'\langle h_1,h_2,\ldots\rangle,\qquad |h_n|=1-2n,
\]
with central $R'$ and differential
\[
 dh_1=s^2,\qquad dh_n=-\sum_{i+j=n}h_ih_j\quad(n\geq2),
\]
where the summation indices are positive.
Define integers $a_1=1$ and $a_n=-\sum_{i+j=n}a_ia_j$.
There are DGA maps fixing $R'$
\[
 R'/(s^2)\xleftarrow{\ \varepsilon\ }P
       \xrightarrow{\ \Phi\ }A(S,c,r),
\]
where
\[
 \varepsilon(h_n)=0,\qquad
 \Phi(h_n)=a_n\tau u^{n-1}\bigl(s+(2n-1)cv\bigr).
\]
The map $\varepsilon$ is always a quasi-isomorphism.
The cohomology map of $\Phi$ in each degree $-2a$ is
$\operatorname{diag}(I_r,cI_r)$.
It is an isomorphism exactly when $c$ is a unit.
\end{proposition}

\begin{proof}
All $h_n$ are odd.
In $d^2h_n$, the scalar terms cancel as
$-s^2h_{n-1}+h_{n-1}s^2$.
Each remaining ordered triple product occurs twice with opposite signs.
Thus $d^2=0$, including in characteristic two.
For $n\geq1$,
\[
 d\Phi(h_n)=a_nu^{n-1}\bigl(s^2+2(n-1)csv\bigr),
\]
because $(s-cv)(s+(2n-1)cv)=s^2+2(n-1)csv$.
For $i+j=n$, one also has
\[
 \Phi(h_i)\Phi(h_j)=a_ia_j u^{n-1}\bigl(s^2+2(n-1)csv\bigr).
\]
The recurrence for $a_n$ proves that $\Phi$ commutes with the differential.
The augmentation $\varepsilon$ does so directly.

Consider the auxiliary DGA
\[
 K=R'\oplus R'h,\qquad |h|=-1,\qquad h^2=0,\qquad dh=s^2.
\]
The exterior relation is imposed also in characteristic two.
The map $\pi:P\to K$ given by $h_1\mapsto h$ and $h_n\mapsto0$ for $n\geq2$
is a DGA map.
Give $h_n$ weight $n$, $h$ weight one, and $R'$ weight zero.
The differential preserves the increasing weight filtration.

At weight $N\geq1$, words $h_{i_1}\cdots h_{i_\ell}$ correspond to
compositions $i_1+\cdots+i_\ell=N$.
Equivalently, they specify the subset of cuts in $\{1,\ldots,N-1\}$.
The associated graded piece is the exterior module
\[
 R'\otimes_S\Lambda_S(e_1,\ldots,e_{N-1}),
 \qquad d_0=-\Bigl(\sum_{j=1}^{N-1}e_j\Bigr)\wedge(-).
\]
A $q$-cut word with coefficient of degree $d$ has degree $d+q+1-2N$.
Splitting the $k$th factor has sign $(-1)^k$,
which is exactly the sign for insertion after $k-1$ existing cuts.
For $N\geq2$, minus removal of $e_1$ contracts this complex:
\[
 \iota_{e_1}(\omega\wedge x)+\omega\wedge\iota_{e_1}(x)=x,
 \qquad \omega=\sum_{j=1}^{N-1}e_j.
\]
This identity uses only the exterior relations and integral signs.
In weights zero and one, $\operatorname{gr}\pi$ is the identity.

The cone of $\pi$ has acyclic associated graded pieces.
Every cone cycle has a greatest weight.
Subtract a differential whose top component is its primitive, and descend in weight.
Finite support makes this process terminate.
Hence $\pi$ is a quasi-isomorphism, without a convergence hypothesis.
Multiplication by the monic element $s^2=t^{2r}$ is injective on $R'$,
so $K\to R'/(s^2)$ is a quasi-isomorphism.
This proves the assertion for $\varepsilon$.

In degree $-2a$, use the source basis
$u^a(1,t,\ldots,t^{2r-1})$ and the target basis
\[
 u^a(1,t,\ldots,t^{r-1},v,vt,\ldots,vt^{r-1}).
\]
The map fixes $t,u$, and $t^{r+i}$ maps to $cvt^i$.
Its matrix is therefore $\operatorname{diag}(I_r,cI_r)$.
Its kernel is $\operatorname{Ann}_S(c)^r$, and its cokernel is $(S/cS)^r$.
A unit $c$ gives the asserted formality zigzag.
Together with the $r=1$ inclusion, this proves sufficiency in
Theorem~\ref{v4-arith:ks:coefficient-formality}.
\end{proof}

\section{The sixth operation on the full centre}
\label{v4-arith:ks:sixth-operation}

The positive generator $z$ changes the obstruction calculation when $2$ is invertible.
It permits a primitive for $m_4$, but that change produces a sixth operation.
A bar cycle containing $z$ detects the resulting class.

Assume throughout this section that $S\ne0$ and $2\in S^\times$.
Write
\[
 H_C=S[u,t,z]/(t^r-2uz,z^2),\qquad |u|=-2,\quad |t|=0,\quad |z|=2.
\]
Its normal basis consists of $u^at^iz^j$ with $a\geq0$, $0\leq i<r$, and $j=0,1$.
Monic division by $t^r-2uz$ gives the same contraction as before.
The transfer lemma, with coefficient algebra $S[u,z]/(z^2)$, gives
\[
 m_4=-uF\smile F,
\]
\[
 m_6(x_1,\ldots,x_6)=u^2\bigl(
 F_{12}F(F_{34},F_{56})+F(F_{12},F_{34})F_{56}\bigr),
 \qquad F_{ij}=F(x_i,x_j).
\]
Indeed, the only nonzero partitions in arity six are $(2,4)$ and $(4,2)$.
Every odd-arity operation vanishes by degree.
In particular the augmentation
$\epsilon:H_C\to S$, $t,u,z\mapsto0$, kills $m_4$ and $m_6$.

To construct a primitive for $m_4$, consider the associative algebra
\[
 S[u,t,z,\alpha,\beta]/(t^r-2uz-\alpha,z^2-\beta),
 \qquad |\alpha|=0,\quad |\beta|=4.
\]
Monic division makes it free over $S[u,\alpha,\beta]$ on $t^iz^j$,
with $0\leq i<r$ and $j=0,1$.
Let $F,G$ be the coefficients of $\alpha,\beta$ in its multiplication,
evaluated at $\alpha=\beta=0$.
They have degrees $0,-4$, respectively.
The first $F$ agrees with the multiplication defect already used.
Associativity in the two parameters gives $\delta F=\delta G=0$.

Define normal-form coefficient derivatives by
\[
 \theta_u(u^at^iz^j)=a u^{a-1}t^iz^j,\qquad
 \theta_z(u^at^iz^j)=j u^at^iz^{j-1}.
\]
The first expression is zero when $a=0$.
These maps have degrees $2,-2$ and satisfy
\begin{equation}\label{v4-arith:ks:derivative-defects}
 \delta\theta_u=-2zF,\qquad
 \delta\theta_z=-2uF+2zG.
\end{equation}
For a direct verification take $x=u^at^iz^j$ and $y=u^bt^\ell z^m$.
Put $A=a+b$, $I=i+\ell$, and $J=j+m$.
The formulas are
\[
 F(x,y)=
 \begin{cases}
 u^At^{I-r}z^J,&I\geq r,\ J\leq1,\\
 0,&\text{otherwise},
 \end{cases}
\]
\[
 G(x,y)=
 \begin{cases}
 u^At^I,&I<r,\ J=2,\\
 2u^{A+1}t^{I-r}z^{J-1},&I\geq r,\ J\geq1,\\
 0,&\text{otherwise}.
 \end{cases}
\]
Substitution proves \eqref{v4-arith:ks:derivative-defects} over $S$.
No derivative of $t^r$ occurs, so no inverse of $r$ is required.

For even internal degrees use the cup product
$(P\smile Q)(\mathbf x,\mathbf y)=P(\mathbf x)Q(\mathbf y)$.
For the two binary cochains put
\[
 (F\circ G)(x,y,z)=F(G(x,y),z)-F(x,G(y,z)).
\]
Expansion of the Hochschild differential gives
\[
 \delta(F\circ G)=G\smile F-F\smile G.
\]
It follows that the ternary cochain of degree $-2$
\begin{equation}\label{v4-arith:ks:fourth-primitive}
 g=\tfrac12(\theta_z\smile F+\theta_u\smile G)-z(F\circ G)
\end{equation}
satisfies $\delta g=-uF\smile F=m_4$.
The two terms with coefficient $z$ cancel by the preceding identity.

Let $Q$ be the tensor-coalgebra coordinate change with components
$q_1=1$, $q_3=-g$, and all other $q_n=0$.
It is invertible because its deviation from the identity lowers tensor length.
The inverse is a finite sum on each finite tensor word.
Conjugating the transferred coderivation gives a minimal structure $m'$ with
\[
 m'_3=m'_4=m'_5=0,
\]
\[
 m'_6=m_6-\sum_{j=0}^3
 m_4(1^{\otimes j}\otimes g\otimes1^{\otimes(3-j)})+g\smile g.
\]
The arity-four identity uses $m'_4=m_4-\delta g$.
The arity-six signs are the morphism signs specified in
Section~\ref{v4-arith:ks:fourth-operation}.
Applying the augmentation gives the decisive identity
\begin{equation}\label{v4-arith:ks:sixth-augmentation}
 \epsilon m'_6=(\epsilon g)\smile(\epsilon g).
\end{equation}
The arity-seven identity makes $m'_6$ a Hochschild cocycle.

\begin{lemma}[a bar cycle detects the sixth operation]
\label{v4-arith:ks:sixth-pairing}
\ClaimStatusProvedHere{}
For $r\geq2$, the cocycle $m'_6$ is not a Hochschild boundary.
\end{lemma}

\begin{proof}
Put $a=t^{r-1}$.
In the reduced bar complex of the augmented algebra $H_C$, use
\[
 b[x_1|\cdots|x_n]
 =\sum_{j=1}^{n-1}(-1)^{j-1}
 [x_1|\cdots|x_jx_{j+1}|\cdots|x_n].
\]
All letters have even internal degree, so their suspensions are odd.
The chains
\[
 \rho=[a|t]-[u|z]-[z|u],\qquad \sigma=[z|z]
\]
are cycles: $b\rho=t^r-2uz=0$ and $b\sigma=z^2=0$.
Write $\mathrm{sh}$ for the signed shuffle product.
It interleaves words while preserving each word's order.
Its sign is the sign of the permutation of the suspended letters.
It commutes with $b$ because products between different words cancel in pairs,
while products inside one word give its own bar differential.
Thus $\gamma=\rho\mathbin{\mathrm{sh}}\rho\mathbin{\mathrm{sh}}\sigma$ is a length-six bar cycle.

On its letters, \eqref{v4-arith:ks:fourth-primitive} gives
\[
 \epsilon g(z,a,t)=\epsilon g(z,t,a)=\tfrac12,
 \qquad \epsilon g(u,z,z)=\tfrac12.
\]
All other triples contribute zero.
For $r=2$ the equality $a=t$ identifies the first two triples.
The complete list of nonzero contributions to
$((\epsilon g)\smile(\epsilon g))(\gamma)$ is
\[
 \begin{array}{c|r|r}
 \text{word}&\text{shuffle coefficient}&\text{cochain value}\\ \hline
 \bigl[z|a|t|z|a|t\bigr]&2&1/4\\
 \bigl[z|a|t|u|z|z\bigr]&-2&1/4\\
 \bigl[u|z|z|z|a|t\bigr]&-2&1/4
 \end{array}
\]
These coefficients are obtained by choosing the two ordered positions
of each of the three two-letter words and summing their permutation signs.
The table remains valid after $a=t$, with identical words collected.
Therefore
\[
 \epsilon m'_6(\gamma)=\tfrac12-\tfrac12-\tfrac12=-\tfrac12.
\]
If $m'_6=\delta K$, its exterior Hochschild terms vanish after $\epsilon$
because every bar letter has zero augmentation.
Its interior terms evaluate $K$ on $b\gamma=0$, up to an overall sign.
Their sum is zero.
The element $-1/2$ is a unit in the nonzero ring $S$, giving the contradiction.
\end{proof}

\begin{theorem}[formality of the full relative centre]
\label{v4-arith:ks:full-centre-formality}
\ClaimStatusProvedHere{}
For every commutative unital ring $S$ and every $r\geq1$,
\[
 C_{S,r}\text{ is formal over }S
 \quad\Longleftrightarrow\quad r=1\ \text{or}\ S=0.
\]
For a nonzero $S$ and $r\geq2$, it also has no formality zigzag
that fixes $S[t]$.
\end{theorem}

\begin{proof}
For $r=1$, the inclusion $S[u,z]/(z^2)\hookrightarrow C_{S,1}$
is a quasi-isomorphism by monic division.
For $S=0$ the statement is immediate.
Suppose $r\geq2$ and $S\ne0$.
If $2$ is a nonunit, a formality zigzag would truncate to one for
$A(S,2,r)$, contradicting Theorem~\ref{v4-arith:ks:coefficient-formality}.

If $2$ is a unit, use the minimal structure $m'$ constructed above.
Every homogeneous component of $H_C$ is free over $S$.
The finite comparison lemma transports a hypothetical formality zigzag
through arity six to a morphism from $m'$ to strict $H_C$.
Normalize its linear component to the identity.
Parity forces its second and fourth components to vanish.
Let $a_3,a_5$ denote its third and fifth components.
The arity-four and arity-six identities give
\[
 \delta a_3=0,\qquad m'_6=\delta a_5+a_3\smile a_3.
\]
The self-insertion
\[
 \begin{split}
 (a_3\circ a_3)(x_1,\ldots,x_5)={}&a_3(a_3(x_1,x_2,x_3),x_4,x_5)\\
 &+a_3(x_1,a_3(x_2,x_3,x_4),x_5)\\
 &+a_3(x_1,x_2,a_3(x_3,x_4,x_5))
 \end{split}
\]
satisfies $\delta(a_3\circ a_3)=-2a_3\smile a_3$.
To verify this identity, expand each differential by the displayed Hochschild rule.
The terms with an adjacent multiplication inside an inner or outer $a_3$
cancel by $\delta a_3=0$.
The two remaining end decompositions each give $-a_3\smile a_3$.
Thus $m'_6=\delta(a_5-\tfrac12a_3\circ a_3)$,
contrary to Lemma~\ref{v4-arith:ks:sixth-pairing}.
Finally, forgetting the extra $S[t]$-action would turn any relative
formality zigzag into an $S$-linear one.
\end{proof}

The full centre therefore remains nonformal over every field when $r\geq2$.
At characteristic two its fourth operation is obstructed.
When two is invertible, the fourth operation is exact and the sixth operation is obstructed.
The nonpositive truncation has the different coefficient criterion of
Theorem~\ref{v4-arith:ks:coefficient-formality}.

\section{Coefficient change and the arithmetic comparison}
\label{v4-arith:ks:coefficient-change}

For a ring map $S\to S'$, the normal bases and monic-division contraction give
\[
 A(S,c,r)\otimes_S S'=A(S',c',r),\qquad
 H^*(A(S,c,r))\otimes_S S'\cong H^*(A(S',c',r)),
\]
where $c'$ is the image of $c$.
The same argument gives
\[
 C_{S,r}\otimes_S S'=C_{S',r},\qquad
 H^*(C_{S,r})\otimes_S S'\cong H^*(C_{S',r}),
\]
with
\[
 H^*(C_{S,r})=S[u,t,z]/(z^2,t^r-2uz).
\]
These are assertions for the displayed complexes and their split normal forms.
They hold without flatness of $S'$.
The relative-centre resolution also extends to $S'[t]$ on its two free generators,
so its multiplicative identification persists after coefficient change.

\begin{corollary}[the exceptional coefficient]
\label{v4-arith:ks:exceptional-two}
\ClaimStatusProvedHere{}
For $r\geq2$, the nonpositive part of the integral relative centre is nonformal.
Its extension from $\mathbb Z$ to $S$ is formal exactly when $2\in S^\times$.
In particular it is formal over $\mathbb Q$ and every field of odd characteristic,
and nonformal over every field of characteristic two and over $\mathbb Z/4$.
The full $C_{S,r}$ is nonformal whenever $2$ is a nonunit in $S$ and $r\geq2$.
For $r=1$, both the full centre and its nonpositive part are formal over every $S$.
\end{corollary}

\begin{proof}
The assertions about the nonpositive part follow from
Proposition~\ref{v4-arith:ks:centre-truncation} and Theorem~\ref{v4-arith:ks:coefficient-formality}.
For any cochain DGA $D$, the complex with $D^n$ in negative degrees,
$Z^0(D)$ in degree zero, and zero in positive degrees is a sub-DGA.
This truncation functor preserves quasi-isomorphisms.
Applying it to a formality zigzag for $C_{S,r}$ would give a formality
zigzag for $A(S,2,r)$.
This proves the stated nonformality of the full centre.
For $r=1$, the inclusion
$S[u,z]/(z^2)\hookrightarrow C_{S,1}$ is a quasi-isomorphism:
monic division identifies cohomology by substituting $t=2uz$.
The analogous inclusion of $B$ proves the nonpositive assertion.
\end{proof}

\begin{corollary}[localization of the full centre]
\label{v4-arith:ks:centre-localization}
\ClaimStatusProvedHere{}
If $2\in S^\times$, then $C_{S,r}[u^{-1}]$ is formal over
$S[t,u,u^{-1}]$ for every $r\geq1$.
\end{corollary}

\begin{proof}
After localization, $z=u^{-1}v$ gives the inverse of
\[
 A(S,2,r)[u^{-1}]\xrightarrow{\ j\ }C_{S,r}[u^{-1}].
\]
The zigzag in Proposition~\ref{v4-arith:ks:positive-zigzag} fixes $S[t,u]$.
Localization at its central even element $u$ is exact on modules,
so it preserves both quasi-isomorphisms.
Its cohomology target is
$S[t,u,u^{-1}]/(t^{2r})$, with $z=t^r/(2u)$.
\end{proof}

\begin{example}[unit and nonunit parameters]
For $S=\mathbb F_2$ and $r=2$, the model $A(S,0,2)$ has
$m_4(t,t,t,t)=-u=u$ and is nonformal.
The model $A(S,1,2)$ is formal by the explicit zigzag.
It is not the reduction of $A(\mathbb Z,2,2)$, whose parameter reduces to zero.
Over $\mathbb Z/4$, the parameters $2$ and $3$ likewise give
nonformal and formal models, respectively.
The criterion is invertibility of $c$, rather than its nonvanishing.
For the zero ring, the unique element is a unit and every displayed algebra is zero.
\end{example}

The global arithmetic object
$\mathcal V^{\mathrm{prim}}$
of Definition~\ref{v4-arith:def:global-iwahori-hecke} is a Hochschild trace class.
The centre above is a relative derived endomorphism algebra.
A comparison between them requires a multiplicative map with a specified base,
source, target, and cohomology isomorphism.
An equality of degree-zero centres does not supply these data.

\begin{proposition}[scope of the centre obstruction]
\label{v4-arith:ks:centre-comparison-scope}
\ClaimStatusProvedHereConditional{}
Let $\mathcal Z$ be an associative $S$-DGA model for an arithmetic derived centre.
Suppose a finite multiplicative $S$-linear quasi-isomorphism zigzag identifies
$\tau_{\leq0}\mathcal Z$ with $A(S,2,r)$.
If $r\geq2$ and $2$ is a nonunit, then $\mathcal Z$ is nonformal over $S$.
If the full $\mathcal Z$ is multiplicatively quasi-isomorphic to $C_{S,r}$,
it is nonformal for every nonzero $S$ and $r\geq2$.
A relative dg Morita equivalence with $E_r$ over $S[t]$, together with
its bimodule comparison, is one possible source of the required identification.
\end{proposition}

\begin{proof}
A formality zigzag for $\mathcal Z$ would truncate to one for
$\tau_{\leq0}\mathcal Z$.
Composing with the hypothesized comparison would make $A(S,2,r)$ formal,
contrary to Theorem~\ref{v4-arith:ks:coefficient-formality}.
The full comparison transfers the classification in
Theorem~\ref{v4-arith:ks:full-centre-formality}.
\end{proof}

The positive zigzag fixes $S[t,u]$ on the nonpositive algebra.
The full centre has the additional sixth-operation obstruction of
Theorem~\ref{v4-arith:ks:full-centre-formality}.
For the global Iwahori category, the required relative bimodule comparison
is an additional construction, independent of rectification of the $E_1$ operad.
